\documentclass[10pt]{beamer}

\usetheme[progressbar=foot]{metropolis}
\usepackage{appendixnumberbeamer}

\usepackage{booktabs}
\usepackage[scale=2]{ccicons}

\usepackage{pgfplots}
\usepgfplotslibrary{dateplot}

\usepackage{xspace}
\usepackage{xcolor}
\usepackage{hyperref}
\hypersetup{colorlinks=true, linkcolor=, urlcolor=blue}

\DeclareMathOperator{\stdev}{stdev}
\DeclareMathOperator{\var}{var}
\DeclareMathOperator{\cov}{cov}
\DeclareMathOperator{\corr}{corr}
\DeclareMathOperator{\prob}{prob}
\DeclareMathOperator{\n}{n}
\DeclareMathOperator{\N}{N}
\DeclareMathOperator{\Cov}{Cov}

\newcommand{\hlf}{\frac{1}{2}}
\newcommand{\bi}{\begin{itemize}}
\newcommand{\ei}{\end{itemize}}
\newcommand{\im}{\item}
\newcommand{\D}{\mathrm{d}}
\newcommand{\E}{\mathrm{e}}
\newcommand{\mye}{\ensuremath{\mathsf{E}}}
\newcommand{\myreal}{\ensuremath{\mathbb{R}}}
\newcommand{\bq}{\begin{equation}}
\newcommand{\eq}{\end{equation}}
\newcommand{\bfr}{\begin{frame}}
\newcommand{\efr}{\end{frame}}
\newcommand{\bas}{\begin{align*}}
\newcommand{\eas}{\end{align*}}

\setbeamertemplate{frame footer}{PCA and Factor Models, BUSI 521/ECON 505}

\title{Principal Components and Factor Models}

\date{}
\author{Kerry Back\\
BUSI 521/ECON 505\\
Rice University\\[1em]
\small\href{https://colab.research.google.com/github/kerryback/busi521/blob/main/FactorAnalysisPrincipalComponents.ipynb}{Open companion notebook in Colab}}


\begin{document}

\maketitle

\section{Statistical Factor Models}

\begin{frame}{Setup}
Let $R_t$ be an $N$--vector of excess returns.  Let $K<N$.

Null Hypothesis:

$$R_t = \mu + Bf_t + e_t$$

for unknown  $\mu \in \mathbb{R}^N$, unknown $B \in \mathbb{R}^{N\times K}$, unobserved $f_t$ with mean zero and covariance matrix equal to the identity, and unobserved $e_t$ with mean zero and unknown diagonal covariance matrix $D$, independent of $f_t$.  
\end{frame}

\begin{frame}{Identification}
The null hypothesis does not identify $B$ or $f_t$, because if $A$ is a $K \times K$ orthogonal matrix, then $Bf_t = BAA'f_t$ and the covariance matrix of $A'f_t$ is also the identity, so $\hat B := BA$ and $\hat f_t := A'f_t$ also satisfy the null hypothesis.

The standard additional restriction made to achieve identification is that $B'D^{-1}B$ is diagonal.

\medskip
This is called \alert{rotational indeterminacy}: an orthogonal matrix preserves lengths and angles, so multiplying by $A$ is a rotation (and/or reflection) of the factor space.  Any rotation of the factors is observationally equivalent.
\end{frame}

\begin{frame}{Inferring the Factors}
Under the null and assuming joint normality, the GLS estimate at each date $t$
$$\hat{f}_t = (B'D^{-1}B)^{-1}B'D^{-1}(R_t - \mu)$$
maximizes the conditional likelihood $p(R_t \mid f_t)$.  Equivalently, it solves
$$\min_f \; \sum_{i=1}^N \frac{(R_{it} - \mu_i - b_i'f)^2}{d_i}$$

\bi
\im Assets whose variation is mostly factor-driven (small $d_i$) receive \alert{more} weight.
\im Assets with large idiosyncratic variance receive less weight.
\ei

\medskip
In python: \texttt{sklearn.decomposition.FactorAnalysis}
\end{frame}

\begin{frame}{GLS Creates Traded Factors}
\begin{itemize}
\item The GLS formula
$$f_t = (\hat{B}'\hat{D}^{-1}\hat{B})^{-1}\hat{B}'\hat{D}^{-1}(R_t - \hat{\mu})$$
shows that each factor is a linear combination of excess returns minus a constant.  The constant is the mean of the factor.

\item Arbitrage Pricing Theory $\Rightarrow$ factor is a pricing factor.

\item For pricing factors, we do not require the mean to be zero.  Furthermore, constants can be dropped without affecting pricing relations.

\item If we drop the constant and scale the regression coefficients to sum to one, then we have a factor that is 100\% invested in risky assets and 100\% short the risk-free asset.
\end{itemize}

\smallskip
{\small \href{https://colab.research.google.com/github/kerryback/busi521/blob/main/FactorAnalysisPrincipalComponents.ipynb}{$\rightarrow$ Open companion notebook in Colab}}
\end{frame}

\section{Principal Components}

\begin{frame}{Eigenvectors of the Covariance Matrix}
\bi
\im Let $\Sigma = \cov(R_t)$ and let $C$ be the matrix of eigenvectors of $\Sigma$, ordered by eigenvalue (largest first).
\im The eigenvectors are orthonormal: $C'C = CC' = I$.
\im Let's focus on the $K$ largest eigenvalues, for some $K$. 
\im Partition $C = (C_K \mid C_{N-K})$ where $C_K$ contains the first $K$ eigenvectors.
\im Notice that
$$C_KC_K' + C_{N-K}C_{N-K}' = \begin{pmatrix} C_K & C_{N-K}\end{pmatrix}\begin{pmatrix}C_K' \\ C_{N-K}'\end{pmatrix} = CC' = I
$$
\ei 
\end{frame}

\begin{frame}{PCA as a Factor Model}
\bi 
\im Define $K$ dimensional $f_t$ and $N-K$ dimensional $e_t$ as
\begin{align*}
f_t &= C_K'(R_t - \hat{\mu}) \\
e_t &= C_{N-K}'(R_t - \hat{\mu})
\end{align*}
\im Then 
\begin{align*}
C_Kf_t + C_{N-K}e_t &= C_KC_K'(R_t - \hat{\mu}) + C_{N-K}C_{N-K}'(R_t - \hat{\mu}) \\
&= (C_KC_K' + C_{N-K}C_{N-K}')(R_t - \hat{\mu}) = R_t - \hat{\mu}
\end{align*}
\im Rearrange, set $B = C_K$, and  $\varepsilon_t = C_{N-K}e_t$ to obtain
$$R_t = \mu + Bf_t + \varepsilon_t $$
\ei
\end{frame}

\begin{frame}{Covariance Structure}
The eigenvector equation is $\Sigma C = C\Lambda$, where $\Lambda$ is the diagonal matrix of eigenvalues.  Post-multiplying by $C'$ gives the eigen-decomposition $\Sigma = C\Lambda C'$.  This implies
\bi
\im $\cov(f_t) = C_K'\Sigma\, C_K = \Lambda_K$ \quad (diagonal matrix of the $K$ largest eigenvalues)
\im $\cov(\varepsilon_t) = C_{N-K}\,C_{N-K}'\Sigma\, C_{N-K}\,C_{N-K}' = C_{N-K}\,\Lambda_{N-K}\,C_{N-K}'$ \quad (\alert{not} diagonal)
\im $\cov(f_t, \varepsilon_t) = C_K'\Sigma\, C_{N-K}\,C_{N-K}' = \Lambda_K C_K'C_{N-K}\,C_{N-K}' = 0$
\ei
The residual covariance matrix is not diagonal, but its eigenvalues are the smallest eigenvalues of $\Sigma$.
\end{frame}

\begin{frame}{How Much Variance Is Explained?}
Each eigenvalue $\lambda_i$ equals the variance of the $i$th principal component.

The fraction of total variance explained by $K$ components is
$$\frac{\lambda_1 + \cdots + \lambda_K}{\lambda_1 + \cdots + \lambda_N}$$

For the 25 size/BM portfolios, the first eigenvalue alone accounts for about 82\% of total variance.  Three components capture about 93\%.

In python: \texttt{sklearn.decomposition.PCA}
\end{frame}


\section{Optimization Perspectives}

\begin{frame}{PCA $=$ Maximum-Variance Portfolios}
The first principal component solves
$$\max_{w}\; w'\Sigma w \quad \text{subject to}\quad w'w = 1$$
The FOC is $\Sigma w = \lambda w$\,---\,the eigenvector equation.

The second component solves the same problem with the additional constraint $w'w_1 = 0$.

Successive components are maximum-variance portfolios orthogonal to all previous ones.
\end{frame}

\begin{frame}{Statistical FA $=$ Minimum Residual-Variance Portfolios}
The factor-analysis FMPs solve
$$\min_{w}\; \hlf w'Dw \quad \text{subject to}\quad B'w = e_i$$

This minimizes idiosyncratic risk while targeting a unit loading on factor $i$ and zero loadings on all other factors.

The FOC gives $w = D^{-1}B(B'D^{-1}B)^{-1}e_i$, and the resulting portfolio returns reproduce exactly the GLS factor estimates.
\end{frame}


\end{document}
